\documentclass[11pt]{article}
\usepackage[T1]{fontenc}
\usepackage{lmodern}
\usepackage{microtype}
\usepackage[margin=1in]{geometry}
\usepackage{amsmath,amssymb,amsthm,mathtools}
\usepackage{algorithm}
\usepackage{algpseudocode}
\usepackage{booktabs}
\usepackage[numbers,sort&compress]{natbib}
\usepackage[hidelinks]{hyperref}
\usepackage[capitalise,noabbrev]{cleveref}
\setlength{\emergencystretch}{3em}

\theoremstyle{plain}
\newtheorem{theorem}{Theorem}[section]
\newtheorem{proposition}[theorem]{Proposition}
\theoremstyle{definition}
\newtheorem{definition}[theorem]{Definition}

\newcommand{\lean}[1]{\path{#1}}
\newcommand{\N}{\mathbb{N}}
\newcommand{\at}[2]{#1|_{#2}}
\newcommand{\btw}[3]{#1|_{(#2,\,#3]}}
\newcommand{\hd}[1]{\lvert #1 \rvert}
\newcommand{\pre}{\preceq}

\title{A Knowledge Store as of a Cut:\\
Append Chains, Consistent Cuts, and What Erasure and Scope Do to a Snapshot}
\author{Alberto Serrano-Calva}
\date{November 10, 2026}

\begin{document}
\maketitle

\begin{abstract}
A query on a knowledge store can ask for the store as of a past moment. In the store studied here
that reading exists for one relation, a citation graph indexed by commit, and for the files of
its repository; every other store is read as it is now. We design the rest. A store is modelled as a family of append chains, a cut as a
prefix length for each, and the read at a cut as the prefixes. Ten laws are machine-checked in
Lean~4 over Mathlib. A read at a cut is unchanged by growth of the store, and a reader has that
property exactly when it is a function of the prefixes. Reads at cuts are ordered and compose with
reads between cuts, so a job's checkpoint is a cut. A derived store kept by generation reads what
recomputation reads, with at most one generation of a key alive at a cut. A narrower grant's
snapshot is the restriction of a wider one's, and positions on a shared chain count what a filter
hides. Nothing erased is visible at any cut when an erasure leaves holes; a hole must keep its
key, or a read by key returns a superseded version; and a seal keeps a retired identity, with
its versions, out of earlier cuts. Every frontier resolves to one greatest consistent cut. The
digest of each prefix certifies a cut, survives an erasure that keeps commitments, and attests
no erasure. A store that drops old generations is exact from its horizon on, never wrong under
it, and unable to tell what it lost. The last law gives the rule for an incomplete history: a
read is taken at the greatest consistent cut among the positions the stores still serve, which
never lies after the cut asked; a chain that nothing reads is lowered alone, and a chain that
others read takes them down with it. Reading the code against the usual picture of ledgers and
derived stores locates the one missing ledger, the past of the map from paths to content, and a
log of about 400 bytes a landing on a synthetic fixture supplies it. Measurements on synthetic
fixtures give the cost of its reads and show what an erasure leaves behind today.
\end{abstract}

\section{Introduction}

The system is a personal knowledge store \cite{serranocalva2026palace}: notes, dialogue, papers and
code land through one writer into a dozen stores (content-addressed raw objects, a catalog from
paths to content, version chains, a chat log, attestations, a vector plane, a lexical index, edge
stores), and a sentence of a query algebra reads them. A sentence may carry
\texttt{as of "<cut>"}. Today it moves one relation. The citation graph, which is
stored by commit, is read at the commit named, and a file of the repository is read as that
commit's blob; every other read is of the present, and the answer says so.

Finishing the adjective is not a matter of adding a timestamp filter. The stores have no common
clock: each has its own append order, and wall time is by rule never an ordering key. Some stores
keep their past and some are rewritten in place. Content can be erased, and an erased item must
not return through a snapshot. The asker's grant is today's, not the cut's. And a history is never
complete: it starts somewhere, and retention may move the start. Each of these changes what
``the store as of a cut'' has to mean, and each can be stated and proved before any code reads a
store at a cut. The contributions:
\begin{enumerate}
\item A model (\Cref{sec:model}) of a multi-store as a family of append chains with cuts, reads at
  and between cuts, derived views, a grant, an erasure and certificates.
\item Ten laws (\Cref{sec:laws}), SNP-1 to SNP-10, machine-checked as 175 theorems, with eighteen
  failure modes as theorems too, among them the reader that ignores the cut, the store rewritten
  in place, a filter applied before the cut, an erasure that deletes rows, a hole that loses its
  key, a digest with collisions, a pruned store read under its horizon, and a record of a cut
  placed on a chain that content reads.
\item A rule for incomplete histories (\Cref{thm:horizon}): among the positions the stores still
  serve there is one greatest consistent cut under the cut asked, a read taken there holds no
  event from after the cut it names, and it differs from the cut asked only on the chains that
  could not serve when no other chain reads them.
\item A design (\Cref{sec:design}) read off the code: the one ledger that is missing, how each
  derived store is a function of a cut, what an erasure must reach, and what an answer says,
  with measurements on synthetic fixtures (\Cref{sec:eval}).
\end{enumerate}

\section{Background and related work}

\paragraph{Consistent cuts.}
Lamport's happened-before relation \cite{lamport1978time} orders events without a global clock; a
global state is a consistent cut, a set of events closed under that order
\cite{chandy1985snapshots}. Vector timestamps characterize the order, and the consistent cuts
form a lattice \cite{mattern1989virtual}; the down-sets of any partial order form a distributive
lattice \cite{birkhoff1937rings,davey2002lattices}. The store's own clock design already takes a
cut to be a down-set of a derived event order with a per-store frontier, and its ledger of laws
already has their closure under meet and join. Here there is one writer, so no marker protocol
is needed; what is needed is the greatest consistent cut under a named frontier, also when a
store cannot serve a position.

\paragraph{Temporal and multiversion data.}
Temporal databases keep rows valid over a period, and system-versioned tables answer a query as
of a system time \cite{snodgrass1999temporal,kulkarni2012sql2011}; multiversion concurrency
control keeps versions so that a reader sees a state as of its start
\cite{reed1978naming,bernstein1983mvcc}. Our generations are system-versioned rows whose period
is counted in positions of an append chain, never in wall time (\Cref{thm:forms}). Incremental
view maintenance \cite{gupta1995views} takes the changes to a view's inputs; the delta between
two cuts is that input here.

\paragraph{Hash chains and commitments.}
A hash chain makes an append history tamper-evident
\cite{merkle1987signature,haber1991timestamp}. A chain over contents cannot survive an erasure of
content, and a plain hash of erased content is a test for it. A salted commitment
\cite{halevi1996commitment} whose salt is destroyed avoids both (\Cref{thm:cert}).

\paragraph{Noninterference.}
A grant restricts what a reader may learn \cite{goguen1982security}. A position on a chain is
itself data: it counts every event before it, hidden or not (\Cref{thm:grant}).

\section{The model}\label{sec:model}

\begin{definition}[Stores, cuts, reads]
A \emph{store} is a family of append chains $s : K \to \mathrm{List}\,\alpha$: for each chain key
the events appended to it, in order. Its \emph{head} is $\hd{s}(k) = \mathrm{length}(s\,k)$.
$s'$ \emph{extends} $s$, written $s \pre s'$, when every chain of $s$ is a prefix of its chain in
$s'$; $s \mathbin{+\!\!+} t$ appends chain by chain. A \emph{cut} is $c : K \to \N$, ordered
pointwise. The \emph{read at} $c$ is $\at{s}{c}(k) = \mathrm{take}\,(c\,k)\,(s\,k)$, and the
\emph{read between} $a$ and $b$ is $\btw{s}{a}{b}(k) = \mathrm{drop}\,(a\,k)\,(\at{s}{b}\,k)$. A
\emph{derived view} at a cut is $D(\at{s}{c})$ for a derivation $D$. A reader $F(s, c)$ is
\emph{append invariant} when $F(s \mathbin{+\!\!+} t, c) = F(s, c)$ for all $t$ and all
$c \le \hd{s}$.
\end{definition}

\begin{definition}[A keyed chain, three ways]
For a chain $l$ of writes $(x, v)$, $\mathrm{state}(l)$ maps each key to its last value: the store
rewritten in place. $\mathrm{state}_c(l) = \mathrm{state}(\at{l}{c})$ is the state recomputed at a
cut. The \emph{generational} store $\mathrm{gens}(l)$ has a row per write: the $i$-th write
$(x, v)$ gives $(x, v, \mathrm{born} = i, \mathrm{died})$ with $\mathrm{died}$ the position of
the next write of $x$, absent while the row is the head; a row is \emph{alive} at $c$ when
$\mathrm{born} \le c < \mathrm{died}$, and $\mathrm{read}(\mathrm{gens}(l), c, x)$ is the value of
the row of $x$ alive at $c$. Retention at a horizon $h$, $\mathrm{prune}_h$, drops the rows with
$\mathrm{died} \le h$.
\end{definition}

\begin{definition}[Grant, erasure, seal]
For a set of chains $G$ and an event predicate $p$, the \emph{snapshot} is
$\mathrm{snap}_{G,p}(s, c) = \mathrm{restrict}_{G,p}(\at{s}{c})$: the chains in $G$, each read at
the cut and filtered by $p$. For an erasure predicate $E$, $\mathrm{redact}_E(s)$ replaces each
erased event by a hole, $\mathrm{visible}_E(s, c)$ is $\at{s}{c}$ without the erased events, and
$\mathrm{compact}_E(s)$ deletes them. For a keyed chain, $\mathrm{holes}_E(l)$ replaces the value
of each erased write by a hole and keeps its key. In a chain of landings $(x, \mathrm{some}\,v)$
and retires $(x, \mathrm{none})$, the \emph{stricter seal} serves at $c$ the value $x$ held at
$c$ unless a retire of $x$ lies after $c$, and as its \emph{history} the landings of $x$ inside
$c$ that no retire follows; the \emph{looser seal} serves the value unless $x$ is not live at the
head.
\end{definition}

\begin{definition}[Consistency, resolution, certificates]
Let $R\,k\,i\,k'\,j$ say that the $i$-th event of chain $k$ reads from the $j$-th of $k'$. A cut
$c$ is \emph{consistent} when $R\,k\,i\,k'\,j$ and $i < c\,k$ imply $j < c\,k'$. It \emph{resolves}
a frontier $f$ when it is the greatest consistent cut under $f$. Given for each chain a set $S_k$
of \emph{served} positions, it resolves $f$ \emph{among} them when it is the greatest consistent
cut under $f$ with $c\,k \in S_k$ for all $k$. For a step $\sigma : H \to \alpha \to H$ and a seed
$h_0$, the \emph{digest} of a chain is its left fold; the \emph{certificate} of a cut is the cut
with the digest of every prefix, and a store \emph{verifies} it when the cut lies inside the
store and its own prefixes have those digests.
\end{definition}

\section{The laws}\label{sec:laws}

Each statement names the declarations of \lean{formal/Palace/Snapshot.lean} that check it.

\begin{theorem}[SNP-1 and SNP-5: append invariance]\label{thm:append}
For $c \le \hd{s}$ and any $t$, $\at{(s \mathbin{+\!\!+} t)}{c} = \at{s}{c}$ and
$\btw{(s \mathbin{+\!\!+} t)}{a}{c} = \btw{s}{a}{c}$. A reader $F$ is append invariant if and only
if there is a derivation $D$ with $F(s, c) = D(\at{s}{c})$ for all $c \le \hd{s}$. The reader
$F(s, c) = D(s)$, which ignores the cut, is not append invariant.
\end{theorem}
\begin{proof}[Sketch]
Taking a prefix of an appended list within the first list's length ignores the second
(\lean{readAt_grow}, \lean{readBetween_grow}). For the equivalence, $s$ is $\at{s}{c}$ grown by
its own suffix, so an append invariant $F$ satisfies $F(s, c) = F(\at{s}{c}, c)$, and
$D(t) = F(t, \hd{t})$ works; the converse is the first part (\lean{appendInvariant_iff_exists_view}).
The control is \lean{Witness.present_not_appendInvariant}.
\end{proof}

\begin{theorem}[SNP-2 to SNP-4: order, composition, the head]\label{thm:compose}
If $c \le d$ then $\at{s}{c} \pre \at{s}{d}$, and $\at{(\at{s}{c})}{d} = \at{s}{c \sqcap d}$. For
$a \le b \le c$, $\at{s}{a} \mathbin{+\!\!+} \btw{s}{a}{b} = \at{s}{b}$ and
$\btw{s}{a}{b} \mathbin{+\!\!+} \btw{s}{b}{c} = \btw{s}{a}{c}$. A fold of the prefix at $b$ is the
fold at $a$ continued over $\btw{s}{a}{b}$. An event read at $b$ was read at $a$ or lies in
$\btw{s}{a}{b}$, and a key the delta does not write has the same state at both cuts. Finally
$\at{s}{\hd{s}} = s$, so every view at the head is the derivation of the present.
\end{theorem}
\begin{proof}[Sketch]
List identities for \texttt{take} and \texttt{drop} (\lean{readAt_mono}, \lean{readAt_readAt},
\lean{readAt_append_readBetween}, \lean{readBetween_append}, \lean{foldAt_readBetween},
\lean{mem_readAt_iff}, \lean{stateAt_eq_of_not_written}, \lean{readAt_head}, \lean{view_head}).
\end{proof}
The composition is what lets a job keep a checkpoint: its state at $b$ is its state at $a$
continued over the delta, so a run and a rebuild agree. For a job that joins two chains, a pair
read at $b$ was read at $a$ or has a member in the delta (\lean{pair_readAt}).

\begin{theorem}[SNP-5: three ways of holding a derived store]\label{thm:forms}
For every chain $l$ of keyed writes, cut $c$ and key $x$,
$\mathrm{read}(\mathrm{gens}(l), c, x) = \mathrm{state}_c(l)(x)$, and both are the last write of
$x$ inside $\at{l}{c}$. At most one row of $x$ is alive at $c$. Growth of the chain, which writes
closing stamps into $\mathrm{gens}$, changes no read at a cut inside it. There is no function $G$
with $\mathrm{state}_c(l) = G(\mathrm{state}(l), c)$ for all $l$ and $c$.
\end{theorem}
\begin{proof}[Sketch]
Induction on the chain from its end: a landing closes the key's open row at its own position and
opens a row there, which leaves every read at an earlier cut and of another key as it was
(\lean{genRead_gens}, \lean{stateAt_eq_lastWrite}, \lean{length_filter_selects_le_one},
\lean{genRead_gens_append}). For the last part two chains with one final state differ at the
first cut (\lean{Witness.stateOf_forgets}).
\end{proof}

\begin{theorem}[SNP-6: the grant]\label{thm:grant}
If $G \subseteq G'$ and $p$ implies $p'$, then
$\mathrm{snap}_{G,p}(s, c) = \mathrm{restrict}_{G,p}(\mathrm{snap}_{G',p'}(s, c))$. A snapshot
depends only on the chains of $G$ and on the cut's positions on them, and it is append invariant.
Restricting to whole chains commutes with the cut. Filtering events does not: there are stores
that agree on every admitted event and differ in their snapshots at one cut. The cut read
(\Cref{thm:horizon}) on a set of chains that read from no chain outside it is a function of that
set alone.
\end{theorem}
\begin{proof}[Sketch]
Filters compose and are applied after the prefix (\lean{snapshot_narrow}, \lean{snapshot_congr},
\lean{snapshot_grow}, \lean{restrict_readAt}). For the control, $[\mathrm{hidden}, a]$ and $[a]$
read at position one give nothing and $a$ (\lean{Witness.filter_before_cut},
\lean{Witness.position_channel}). For the last part, a resolved cut restricted to a closed set
and zero elsewhere is consistent and served (\lean{ResolvesIn.eq_on_of_closed}; the control is
\lean{Witness.resolution_reads_unnamed}).
\end{proof}
So what must be hidden in number is hidden by chain: where clearances separate rows, a chain per
class. The law is about the rows a grant is shown. It passes to an answer that filters those
rows, and not to a score, whose collection statistics are the reader's own.

\begin{theorem}[SNP-7: erasure]\label{thm:erase}
No erased event is in $\mathrm{visible}_E(s, c)$; a larger erasure shows a sublist; reading the
store with holes commutes with the cut and keeps the head; and if
$\mathrm{redact}_E(s) = \mathrm{redact}_E(s')$ then $D(\mathrm{visible}_E(s, c)) =
D(\mathrm{visible}_E(s', c))$ for every $D$ and $c$. After $\mathrm{compact}_E$ a read at a cut
can hold an event from after it. On $\mathrm{holes}_E(l)$ the state of a key at a cut is its
last write at or under the cut, a hole when that write is erased, and the generational read and
the seek by key agree with it; if a hole loses its key, a seek returns a superseded version, or
one from after the cut. Under the stricter seal, once a chain retires a key no cut before the
retire serves it, and for any growth $t$ a value served by $l \mathbin{+\!\!+} t$ at a cut inside
$l$ was served by $l$; the history served ends at the value served, and at every cut after a
retire it holds only landings after that retire. The looser seal serves everything the stricter
one serves, and a landing after a retire brings the key's past back under it.
\end{theorem}
\begin{proof}[Sketch]
Filtering and mapping commute with \texttt{take} (\lean{not_erased_of_mem_visible},
\lean{redact_readAt}, \lean{erased_invisible}, \lean{stateAt_holes}). The seal is a predicate
on the suffix past the cut, which only grows (\lean{servedAt_of_append},
\lean{historyAt_of_retire_inside}, \lean{servedNow_of_servedAt}). The controls are
\lean{Witness.compact_moves}, \lean{Witness.keyless_hole} and \lean{Witness.servedNow_restores}.
\end{proof}

\begin{theorem}[SNP-8: resolution]\label{thm:lattice}
A cut is consistent exactly when the events under its frontier are a down-set of the order
generated by the chains and by $R$. Every frontier is resolved by exactly one cut, monotonically
in the frontier. If growth only adds references that start at or past the old head, a consistent
cut inside the old head stays consistent. If every source of an event lies at or under a
recorded one on its chain, a cut consistent for the recorded references is consistent for $R$.
Withholding from a cut every row with a source outside it or withheld leaves a set that holds
the sources of its rows, the greatest such, and the whole prefix when the cut is consistent.
\end{theorem}
\begin{proof}[Sketch]
The down-set form is the clock's cut criterion, no generator edge crossing
(\lean{consistent_iff_isCut}). Closure under pointwise meet and join then is the clock's law for
frontiers, already in the store's ledger, and is derived from it, not claimed here
(\lean{Consistent.inf}, \lean{Consistent.sup}). For resolution take on each chain the supremum
of the positions reached by consistent cuts under $f$: it is attained, and a source of an event
under it is under the witness's position on the source's chain (\lean{exists_resolves},
\lean{Resolves.unique}, \lean{Resolves.mono}, \lean{Consistent.of_grown},
\lean{Consistent.of_bound}, \lean{Kept.source}, \lean{kept_of_closed}).
\end{proof}

\begin{theorem}[SNP-9: certificates]\label{thm:cert}
The certificate of a cut of $s$ verifies on every $s' \succeq s$. If the step is collision free,
two stores that verify one certificate have the same read at its cut, hence the same derived
views. The digest at $b$ is the digest at $a$ continued over $\btw{s}{a}{b}$, and the digests at
$c \sqcap d$ and $c \sqcup d$ are, chain by chain, those of $c$ or of $d$. If the chain digests
commitments kept beside the events, a store verifies a certificate after an erasure of contents
exactly when it did before, and the contents then read as the store with holes; if moreover the
commitment is injective, two such stores that verify one certificate were made from the same
events at its cut.
\end{theorem}
\begin{proof}[Sketch]
Append invariance of the prefix; induction from the end for injectivity of the fold
(\lean{verifies_certify_of_extends}, \lean{Verifies.readAt_eq}, \lean{Verifies.view_eq},
\lean{certify_digest_readBetween}, \lean{certify_inf}, \lean{certify_sup},
\lean{verifies_purge_iff}, \lean{contents_purge_committed}); the last part is
\lean{Verifies.readAt_eq_of_committed}. Under addition two chains of one length collide
(\lean{Witness.digest_collides}).
\end{proof}
A certificate thus pins what it digests and nothing else. It attests no erasure, so a copy taken
before an erasure verifies it and still holds what was erased; no seal, since a retire written
after the cut changes what is served at it (\lean{Witness.seal_not_certified}); and no grant.
Its chains are named, and a renamed chain no longer verifies
(\lean{Witness.rename_breaks_certificate}).

\begin{theorem}[SNP-10: the horizon]\label{thm:horizon}
For $h \le c$, $\mathrm{read}(\mathrm{prune}_h(\mathrm{gens}(l)), c, x) = \mathrm{state}_c(l)(x)$.
For any $c$, a value the pruned store returns at $c$ is $\mathrm{state}_c(l)(x)$; but there are
$l$ and $c < h$ at which it returns nothing where a value stood.
Two chains can have one pruned store and differ, at a cut under $h$, in whether a key had been
written. $\mathrm{prune}_h \circ \mathrm{prune}_{h'} = \mathrm{prune}_{\max(h, h')}$. If every
$S_k$ contains $0$, every frontier $f$ is resolved among the served positions by exactly one cut;
that cut lies under the resolution of $f$, equals it when every position is served, and rises
with $f$ and with served sets that grow. If $f$ is consistent, the chains in a set $U$ are read
only by themselves and by chains of a set $V$ that only themselves read, and every other chain
serves its position, then the cut is $f$ with each chain of $U$ at the greatest position it
serves at or under $f$, and each chain of $V$ at the greatest it serves whose events read $U$
inside those positions. Without the hypothesis on $U$ that fails: one chain read low can take
every other chain to its empty prefix.
\end{theorem}
\begin{proof}[Sketch]
A row alive at $c \ge h$ has no closing stamp at or under $h$, so pruning removes no row a read
at $c$ selects; under $h$ a selected row is still the unique row alive in the full store
(\lean{genRead_prune_gens}, \lean{stateAt_of_genRead_prune}, \lean{prune_prune}); the two
controls are \lean{Witness.prune_forgets} and \lean{Witness.prune_indistinct}. Resolution among
served
positions is the argument of \Cref{thm:lattice} with the supremum taken over served consistent
cuts: the supremum on a chain is attained by one of them, so it is served
(\lean{exists_resolvesIn}, \lean{ResolvesIn.unique}, \lean{ResolvesIn.mono},
\lean{ResolvesIn.le_resolves}, \lean{resolvesIn_univ}). Lowering an unread chain breaks no
reference, and no served consistent cut under $f$ is higher (\lean{resolvesIn_of_unread},
\lean{resolvesIn_of_followed}). The controls are \lean{Witness.record_drags},
\lean{Witness.record_follows} and, for sets that shrink, \lean{Witness.served_set_shrinks}.
\end{proof}

\section{The design}\label{sec:design}

\paragraph{The missing ledger.}
The usual picture has ledgers that only append and derived stores that are rewritten and can be
recomputed. Read against the code it fails at one place that matters. The catalog, the map from a
path to the content it holds and the idempotence key of every landing, is rewritten in place: an
upsert, a tombstone flag, a delete. Version chains exist for one of four places; elsewhere the
attestation ledger names a digest and no path, once per action and digest; so for most items the
sequence of contents a path held is recorded nowhere. The vector plane and the lexical index, which
the picture calls rewritten, keep every version's rows and move a flag. They are generations
without a clock. A clock on their rows would be wrong, since a path that returns to earlier bytes
makes the same rows current again, so a version's rows are valid on a set of intervals. The fix
is one ledger: a \emph{landing log} in the catalog's own file, hash chains that each hold one
provenance class, a row per catalog write (chain, position, kind, a keyed hash of the path, path,
identity, digest, a salt, a commitment, the chain hash, the position of the path's next event),
written by the catalog's own methods in the transaction of the row they write. A chain has an id
and its class is a label, so relabelling a class moves no row and breaks no certificate. The
catalog is then the fold of the log (\Cref{thm:forms}), and every version-keyed store has a
past: the log names the version at a cut.

\paragraph{Cuts, names, certificates.}
A cut is one position per chain, about a dozen integers: the landing chains, the chat log, the
attestations, the repository's commit, the generations of model output. It is not a position per
document, which is millions of pairs on a large instance. A cut is recorded, with the digest of every
hashed prefix, when the writer's queue and its handoff are empty, and attested. The records form
a chain of their own that no other chain reads: every landing reads the attestation ledger, so
a record placed there can, by \Cref{thm:horizon}, take a landing chain to its empty prefix with
any chain read below a position it certified. The adjective takes \texttt{head}, a record's id, a commit
(the first record that holds it, with the repository lowered to the commit and resolved), or a
date, which selects the last record whose wall bookmark is not later. Wall time is a lookup key
there and orders nothing. The resolver follows three references, each through a recorded bound
that is sound because there is one writer (\Cref{thm:lattice}): a landing's to the attestation
ledger, an attestation's to its commit, a record's to its frontier. A pass's attestation and a
row of model output read landings too; resolved, they would lower the ledger, and that can take
another landing chain down (\lean{Witness.pass_drags}), so the filter of \Cref{thm:lattice} withholds them
from the read instead.

\begin{algorithm}[t]
\caption{The cut a read under \texttt{as of} is taken at. $S_k$: the positions chain $k$ serves.
$\beta_k(i)$: for the $i$-th event of chain $k$, the heads of the chains it reads, as recorded at
its writing; they rise with $i$.}\label{alg:resolve}
\begin{algorithmic}[1]
\Require a frontier $f$ from the name
\Ensure the greatest consistent cut under $f$ every position of which is served
\State $c \gets f$
\Repeat
  \For{each chain $k$}
    \State $c_k \gets \max\{p \in S_k : p \le c_k\}$ \Comment{$0 \in S_k$}
    \If{some event of $k$ under $c_k$ has $\beta_k(i)(k') > c_{k'}$ for some $k'$}
      \State $c_k \gets$ the least such $i$ \Comment{a binary search}
    \EndIf
  \EndFor
\Until{nothing moved}
\State \Return $c$
\end{algorithmic}
\end{algorithm}

\paragraph{Derived stores.}
By \Cref{thm:append} a derived store can be read at a cut exactly when it is a function of the
prefixes, and by \Cref{thm:forms} the ways of holding it agree, so the choice is cost. Vectors
are a function of a chunk's text and the embedder's identity and are content addressed; the
plane keeps superseded versions, so the vectors at a cut are those of the versions the log names,
at no new bytes. The lexical index is the same, and its collection statistics
\cite{robertson2009bm25} at a cut are those of the units alive at the cut: a cut, like a
clearance, is a mask whose statistics must be its own. The causal edges from session transcripts
to commits and files are recomputed from the transcript as retained at the cut, which is cheap
and also removes a dependence on job order that the stored rows have today. The relations a note
declares in its own text are recomputed too. A commit's citation edges are stamped with the
projection that wrote them, since a later extractor adds edges under an old commit and the
reading would otherwise grow under a cut that did not move. Model output cannot be recomputed
and is kept by generation, the replaced row archived with its interval; a kept row is served
only while every content it was derived from is alive at the cut and not erased. One storage form is
ruled out as the carrier of a snapshot: the vector table's own versioned manifests. An erased
text stays readable through an earlier manifest until a compaction that removes every earlier
manifest (\Cref{sec:eval}), so a manifest version cannot be both the position of a snapshot and
within reach of an erasure. The cut is a predicate on rows taken from the log.

\paragraph{Erasure and scope.}
Both are parameters of a read, not functions of the cut. A purge empties the log's row and keeps
its position, its commitment, its hash and the keyed hash of its path, so that a read by path
finds the hole (\Cref{thm:erase,thm:cert}); the salt goes with the content, so the commitment
cannot be tested against a guess. A retire seals the identity's earlier cuts under the stricter
seal, lifted only by an owner's act. A new note written at a deleted note's path inherits its
identity by default; the looser seal would bring the deleted versions back at earlier cuts, and
without the seal on the history they would stand under the new note at every later one. The grant is the asker's now, an event filter is applied after the read,
and an answer names the cut only on the chains the grant names (\Cref{thm:grant}).

\paragraph{What an answer says.}
A store may be unable to serve the cut asked: it keeps no past yet, or the cut is under its
horizon. The read is then taken at the cut \Cref{alg:resolve} returns, which by
\Cref{thm:horizon} exists, is unique, never lies after the cut asked and rises as stores gain a
past. A chain that can fail to serve a recorded cut is read, by resolved references, only by
the chain of records, so only those chains and the records' are lowered, and the rows that read
them are withheld. The answer is exact for that cut and names it, chain by chain, with
the reason wherever it differs from the cut asked. Reading such a store at the present instead, which is what the
executor does today, is the reader \Cref{thm:append} refutes. At the head the read is the
present (\Cref{thm:compose}), so \texttt{as of "head"} is the sentence without the adjective.

\paragraph{The horizon and retention.}
A chain's horizon is set by the genesis of its log, by retention, or by a reset. Retention is
counted in positions and has two levels. The derived rows of old versions can be dropped without
moving a horizon: they are cache entries of a function of content, a cut that asks for one finds
it pending, and it is derived again from the raw object. Moving the horizon drops raw objects
and empties the log's rows closed under it, which only the store's owner declares. That pass
appends a horizon row first, so a crash leaves rows that are unserved, never a position that is
served and incomplete, and it keeps every commitment, so the chain verifies from its first row. A window counted per identity (the last $n$
versions of each) would leave no horizon to record.

\section{Evaluation}\label{sec:eval}

\paragraph{What is proved.}
\Cref{thm:append,thm:compose,thm:forms,thm:grant,thm:erase,thm:lattice,thm:cert,thm:horizon} are
machine-checked in \lean{formal/Palace/Snapshot.lean} over \lean{formal/Palace/Snapshot/Defs.lean}
on Lean~4 \cite{moura2021lean4}, v4.36.0-rc2, and the Mathlib \cite{mathlib2020} it pins: 175
theorems, eighteen of them the controls, no
\lean{sorry}, and the axioms of all 5{,}213 library declarations audited (\lean{propext},
\lean{Classical.choice}, \lean{Quot.sound} only). The library's export, 56{,}188 declarations
from \lean{Init} up, is replayed by Lean's kernel and five independent checkers. These are
design theorems: no code relies on them yet, so they are not among the 328 statements the
library restates separately and judges against its ledger of laws.

\paragraph{The attacks.}
The definitions compute, so they are attacked as they are with property-based testing
(\lean{formal/PalaceTest/Snapshot.lean}, seeds 430 to 520, 500 draws each): stores of two chains
over three events, keyed chains over three keys, chains of landings and retires, grants and
erasures from bit masks, reads-from relations as edge lists, a base-four digest, served positions
from ten bits. Forty-nine attacks survive every draw. \Cref{alg:resolve}, which the library
does not define, is run as an equivalent alternation, every chain floored and then one pass of
the descent, and agrees on every draw with the pointwise maximum of every consistent served cut
under the frontier. Forty-two controls are refuted within 61 draws: among them a read at a cut past the head, generations
without the closing stamp, an event filter before the cut, an erasure that deletes rows, a seek
over rows that lost their key, an additive digest, the pruned store read under its horizon, a chain that others read lowered alone, and the floor and the
descent each run once, in either order.

\paragraph{What is measured.}
All on synthetic fixtures built by each probe, on one machine (Apple M2 Pro, SQLite 3.54.1,
LanceDB 0.33.0), with other work running; no instance was read. The probes are not in the
repository; the design note describes each fixture, and each build stage commits its harness. \Cref{tab:measure} gives the figures.

\begin{table}[t]
\centering
\small
\begin{tabular}{@{}p{0.34\linewidth}p{0.60\linewidth}@{}}
\toprule
probe & result \\
\midrule
landing log, $10^6$ events over $10^5$ paths, keyed and hashed, one transaction &
  396 to 434 bytes a row with its three indexes; 35{,}000 to 42{,}000 rows a second \\
one path at the middle cut; its seven versions & 3 $\mu$s and 9 $\mu$s through the key index;
  49 to 57 ms when the planner takes the primary key; after a version's row is emptied the
  seek returns the hole \\
the delta of 1{,}000 positions & 0.64 to 0.69 ms \\
every head at a cut & 147 to 166 ms at the middle; 240 to 275 ms at the head \\
the chain's hash from the first row; the stored hash at a position & 0.7 to 0.8 s; 3 $\mu$s \\
lexical index, 500 items, 4 chunks, 3 versions &
  4{,}179 bytes a unit at 122 postings; 0.67 of the units superseded \\
a purge of both versions of a note, on the test suite's throwaway store &
  the words remain in the lexical postings, and in the vector table read at 8 and at 5 of its 10
  versions until a compaction that leaves 1; the version rows and attestations keep the
  digests; the plane's per-path log gives position 1 again \\
the same transcripts and commits, two job orders & 1 causal edge against 0 \\
re-deriving a session of 2{,}000 records & 6 ms \\
\bottomrule
\end{tabular}
\caption{Measurements on synthetic fixtures. Ranges are across two runs.}\label{tab:measure}
\end{table}

The log's cost is small against what the stores already keep: the vector plane and the lexical
index drop no superseded version today, so most bytes of a past are paid already. Reading every
head from the log alone passes a tenth of a second at about a quarter of a million events in a
chain, past which heads are kept at recorded cuts. The purge row orders the build: the stores
that would be read at a cut are the ones an erasure does not reach, so the erasure's reach lands
before the first stage that reads a superseded version.

\section{Limitations}

The model has one writer: a chain is a list, and nothing is said of two writers. Collision
freedom and a binding commitment are hypotheses of \Cref{thm:cert}, not theorems about a hash
function. \Cref{alg:resolve} is attacked, not proved; the library proves the cut it must return
exists and is unique. The laws of the grant and the erasure are laws of a reader that reads
through them; that every reader in the code does is an obligation of the build, and no single
reader in the model applies the resolution, the grant, the erasure and the seal together. A
repository history with merges is a partial order, which the clock's down-set laws cover and
the chain model does not. Every measurement is on a synthetic fixture and cannot yet be repeated from the
repository: the cost of the lexical statistics from a mask at a corpus's size, of a vector
search joined to the heads at a cut, and of a retention pass are not measured. The attestation
ledger still names an erased digest; the note proposes a salted commitment in its place and
leaves the choice to the store's owner.

\section{Future work}

The build is eight stages, each with its tests and its rollback: the landing log;
recorded cuts, names and the resolver; the erasure's reach and the seal; the notes, vectors and
lexical index at a cut; the causal edges at a cut; model output by generation; retention; and
the delta between two cuts. Each stage property-tests its transcription against the Lean
statements, and its laws then enter the ledger.

\section{Conclusion}

A snapshot of a store with many stores is a prefix of each chain at a consistent cut, read
through the present's grant and the present's erasure. Stated that way its properties are
theorems, and the same statements located what the code lacked, a ledger of the catalog's past,
and what it must not do: answer a question about the past from the present.

\bibliographystyle{plainnat}
\bibliography{refs}
\end{document}
